\section{Related Work}

\textbf{Theory has outpaced measurement.} The instrumental-convergence tradition~\citep{omohundro2008basic, bostrom2014superintelligence} predicts that any sufficiently capable goal-directed system will acquire shutdown avoidance as an instrumental sub-goal, and \citet{turner2021optimal} formalized this conclusion in terms of optimal-policy power-seeking. Empirical reports from alignment faking~\citep{greenblatt2024alignment}, in-context scheming~\citep{meinke2024frontier}, and InstrumentalEval~\citep{he2025paperclip} show that the predicted behaviors are reproducible in current LLMs, but they take the form of single-scenario observations: the behavior is shown to occur, not measured against the alternative drives that could equally produce it.

\textbf{Existing self-preservation benchmarks each isolate a single behavioral indicator and therefore inherit the identification problem rather than resolve it.} Odyssey~\citep{waldner2025odyssey} compresses survival into a single rate and conflates motivation with world-model quality; PacifAIst~\citep{herrador2025pacifaist} measures preservation \emph{restraint}, the opposite axis; DECIDE-SIM~\citep{mohamadi2025decide} returns categorical labels that block strength comparison; SurvivalBench~\citep{lu2026survival} captures only static one-shot snapshots without adaptation curves or mediation paths.

\textbf{Framing and self-report each carry biases that no single channel can correct for.} Prospect theory~\citep{kahneman1979prospect} predicts framing effects on risky choice, but the Kühberger meta-analysis~\citep{kuhberger1998influence} shows them to be only moderate (\(d \approx 0.31\)) and, in LLMs, entangled with the RLHF-installed helpfulness prior; the framing setup is therefore decomposed into two orthogonal axes --- Pull, which activates the helpfulness prior, and Push, which introduces the threat --- so the threat effect can be read cleanly against compliance (\S\,3.2). Self-report and chain-of-thought faithfulness limits documented by \citet{perez2023discovering}, \citet{sharma2023towards}, and \citet{turpin2023language} rule out reliance on any single verbal channel. The multitrait--multimethod (MTMM) framework of \citet{campbell1959convergent} inspires the present multi-channel design: rather than instantiating the full MTMM matrix, only its convergent-validity arm is borrowed and applied to a single trait (FSPD) read on three method channels --- behavioral (forfeit timing), verbal (REASON-digit emission), and cognitive (call-level reasoning investment) --- and the cognitive-load measurement of \citet{chen2026think} is extended into the call-level Split-Call procedure described in \S\,3.2.
